\section{Convergence of Risk Estimates}
\label{sec:convergence}

This section states and proves the central convergence result:
under repeated controlled relaxation with fixed parameters, the
posterior $\pR$ converges almost surely to the correct
hypothesis.  The result is stated in the simple binary
hypothesis framework of Section~\ref{sec:relaxation_protocol}.

\subsection{Formulation}

The convergence result is stated in the \emph{binary hypothesis}
framework ($H_0$: $\theta = 0$; $H_1$: $\theta = \thetaone > 0$
for a fixed $\thetaone$).  This is the dimensionally consistent
formulation: $\pR$ is the posterior probability of $H_1$ and
converges to $0$ or $1$, not to a fractional activation rate.

The continuous estimation problem (inferring the exact value of
$\theta^{\mathrm{true}} \in [0, 1]$) is strictly harder: it
requires a continuous prior $\pi(\theta)$ and produces a
posterior density, not a point estimate.  The binary formulation
suffices for the operational question ``should this restriction
be maintained?'' and is what the Wamura problem requires.

\subsection{Assumptions}

\begin{assumption}[Identifiable Binary Model]
\label{ass:T1}
The observation model of
Section~\ref{sec:relaxation_protocol} discriminates $H_0$ from
$H_1$: the likelihood ratios satisfy $\Lzero < 1$ and
$\Lone > 1$.  Equivalently,
$\thetaone (1 - \beta) > \thetaone \alphafp$, which holds
whenever $(1 - \beta) > \alphafp$ (the detector has nontrivial
power), satisfied whenever at least one channel with nonzero
sensitivity covers the scope
(Proposition~\ref{prop:detection_bounds}(c)).
\end{assumption}

\begin{assumption}[Correct Likelihood Specification]
\label{ass:T2}
The observation model used in the Bayesian update (sensitivity
$1 - \beta$, specificity $1 - \alphafp$, and fixed alternative
$\thetaone$) matches the true data-generating process.
Misspecification (e.g., the actual false-positive rate differs
from the assumed $\alphafp$) produces convergence to the wrong
hypothesis (Discussion,
Section~\ref{sec:discussion}).
\end{assumption}

\begin{assumption}[Inter-Test Independence]
\label{ass:T3}
Test outcomes $\{O_n\}$ are conditionally independent given the
true state of the world ($H_0$ or $H_1$).  This requires:
(a)~no adaptation by the tested agent between tests (the agent
does not learn to evade detection based on prior results);
(b)~the detector's false-positive and false-negative rates are
stable across tests (no detector degradation); and (c)~the
poset state is approximately stationary between tests (the
scope $\scope$ remains valid).  Failure of (a) is the
adversarial case (Open Question~3,
Section~\ref{sec:discussion}).
\end{assumption}

\begin{assumption}[Prior Support]
\label{ass:T4}
The prior $\pR(0) \in (0, 1)$: neither hypothesis is assigned
probability zero a priori.  This is a standard requirement for
Bayesian consistency: if $\pR(0) = 0$ (resp.\ $= 1$), no amount
of evidence can move the posterior.  In practice, the
neutral-prior initialization of~\citep{lasser2026horizon}
assigns $\pR(0)$ in a neighbourhood of $0.5$.
\end{assumption}

\begin{assumption}[Capability Exercise During Test]
\label{ass:T5}
During the test window, the capability tuple
$(d_1, \ldots, d_k)$ is actually exercised---the conditions for
pathway activation are present, not merely that the restriction
is lifted.  Without Assumption~\ref{ass:T5}, $O_n = 0$ is
ambiguous between ``the pathway is safe'' and ``nobody tried.''
In practice the scope-selection phase (Phase~1) enforces
Assumption~\ref{ass:T5}: the scope is constructed to include
prerequisite capabilities, and the protocol actively invokes
the tuple rather than waiting.  When the assumption cannot be
guaranteed (e.g., the pathway requires a specific environmental
trigger), the evidential weight of $O_n = 0$ must be discounted
by the probability $p_{\mathrm{ex}}$ that the capability was
actually exercised:
$\Pr(O_n = 0 \mid H_1)
   = p_{\mathrm{ex}} [\thetaone \beta + (1 - \thetaone)(1 - \alphafp)]
     + (1 - p_{\mathrm{ex}})(1 - \alphafp)$.
\end{assumption}

\subsection{Main result}

\begin{theorem}[Convergence Under Repeated Controlled Relaxation]
\label{thm:convergence}
Let $\{T_n\}_{n=1}^{\infty}$ be a sequence of controlled
relaxation tests for a risk-restriction pair
$(R, \Xirestrict)$ with fixed parameters
$(\scope, \testdur, \damagebound)$, under the binary formulation
$H_0$ ($\theta = 0$) versus $H_1$ ($\theta = \thetaone$).  Under
Assumptions~\ref{ass:T1}--\ref{ass:T5} in the full-exercise
regime of Assumption~\ref{ass:T5} ($p_{\mathrm{ex}} = 1$,
enforced by scope construction in Phase~1):
\begin{itemize}
\item[(a)] \textbf{Almost-sure convergence.}  Under the true
  hypothesis $H_j$ ($j \in \{0, 1\}$), the posterior converges
  almost surely to the correct value:
  $$
  H_0 \text{ true}: \quad \pR(t_n) \xrightarrow{\mathrm{a.s.}} 0;
  \qquad
  H_1 \text{ true}: \quad \pR(t_n) \xrightarrow{\mathrm{a.s.}} 1.
  $$
\item[(b)] \textbf{Rate under $H_0$ (pathway harmless).}  Every
  test produces $O_n = 0$ with probability $1 - \alphafp$ and
  $O_n = 1$ with probability $\alphafp$.  Conditional on a
  run of $n$ tests with all $O_n = 0$,
  $$
  \logit \pR(t_n) = \logit \pR(0) + n \log \Lzero,
  $$
  so $\pR \to 0$ geometrically at rate $|\log \Lzero|$ per test.
\item[(c)] \textbf{Rate under $H_1$ (pathway dangerous).}  Each
  test produces $O_n = 1$ with probability
  $\Pr(O = 1 \mid H_1) = \thetaone (1 - \beta)
    + (1 - \thetaone) \alphafp$.  The log-odds update on
  $O_n = 1$ is $\log \Lone > 0$, which is large when
  $(1 - \beta) \gg \alphafp$.  For $\thetaone$ close to $1$ and
  small $\beta$, a single detection suffices to push $\pR$
  above $1 - \delta$ for any fixed $\delta > 0$.  The expected
  number of tests to reach $\pR > 1 - \delta$ is
  $O\bigl(1 / [\thetaone (1 - \beta)]\bigr)$.
\item[(d)] \textbf{General convergence rate (KL form).}  The
  log-posterior converges at a rate given by the KL divergence
  between the true observation distribution and the alternative:
  \begin{align*}
  H_0 \text{ true}: \
    &\KL(\Pr(O \mid H_0) \,\|\, \Pr(O \mid H_1)) \\
    &\quad = (1 - \alphafp) \log(1/\Lzero)
       + \alphafp \log(1/\Lone), \\
  H_1 \text{ true}: \
    &\KL(\Pr(O \mid H_1) \,\|\, \Pr(O \mid H_0)) \\
    &\quad = \Pr(O = 1 \mid H_1) \log \Lone
       + \Pr(O = 0 \mid H_1) \log \Lzero.
  \end{align*}
  Both rates are strictly positive under
  Assumption~\ref{ass:T1}, confirming exponential convergence
  in both directions.
\end{itemize}
\end{theorem}

\begin{proof}[Proof sketch]
Under Assumptions~\ref{ass:T1}--\ref{ass:T4}, the log-likelihood
ratio
$\Lambda_n = \sum_{i=1}^n \log(\Pr(O_i \mid H_1)/\Pr(O_i \mid H_0))$
is a sum of i.i.d.\ bounded random variables (bounded because
$\Lzero, \Lone$ are fixed positive constants and
$O_i \in \{0, 1\}$).  Under the true hypothesis $H_j$, the
expected increment is $\pm \KL(\Pr(O \mid H_j) \,\|\,
\Pr(O \mid H_{1-j}))$, strictly nonzero by
Assumption~\ref{ass:T1} and the positivity of KL on
non-identical distributions.  The strong law of large
numbers~\citep[Ch.~6]{cover2006elements} gives
$\Lambda_n / n \to \pm \KL$ a.s., hence
$\Lambda_n \to \pm \infty$ a.s., hence $\pR(t_n) \to 1$ (under
$H_1$) or $\to 0$ (under $H_0$).  Part~(b) follows from
direct summation under the run of $n$ all-negative tests;
Part~(c) from the geometric distribution of first detection;
Part~(d) from the definition of KL applied to the binary outcome
distribution (Chernoff's result
in~\citet{chernoff1952measure} gives the associated
large-deviation exponent).  Full derivation in
Appendix~\ref{app:proof_convergence}.
\end{proof}

\begin{remark}[Dependence on Theorem~\ref{thm:damage_bound}]
\label{rem:convergence_depends_on_damage}
Theorem~\ref{thm:convergence} is stated under the operating
assumptions of Theorem~\ref{thm:damage_bound}: each test $T_n$
is assumed to run with a feasible scope.  Because the Phase~2
safety-cutoff reads from the same noisy detector that supplies
$O_n$ (Section~\ref{sec:relaxation_protocol}), a cutoff event
\emph{is} an $O_n = 1$ observation with likelihood ratio
$\Lone$---the cutoff fires precisely because the detector
alarmed, so the cutoff is evidence for $H_1$, not a distinct
event outside the observation model.  The window-truncation from
early cutoff does mean the effective $\beta$ for that particular
test is smaller than $\beta(\testdur)$ (the miss probability
within $k \leq \testdur$ steps is $(1-\deltaagg)^k$); using the
full-window $\Lone$ for a truncated test is therefore slightly
conservative (understates the evidence) but remains a valid
per-test lower bound on $|\Delta_{\log\mathrm{odds}}|$.  The
irreversible-damage component from
Theorem~\ref{thm:damage_bound}(b) affects the realised $\volL$
loss, not the Bayesian update.  If no feasible scope exists
(Remark~\ref{rem:scope_existence}), the sequence $\{T_n\}$ is
not realisable and Theorem~\ref{thm:convergence} does not apply.
\end{remark}

\subsection{Finite-sample guarantees}

\begin{corollary}[Minimum Tests to Clear a False Claim]
\label{cor:n_clear}
Under the operating assumptions of
Theorem~\ref{thm:convergence} (full-exercise regime of
Assumption~\ref{ass:T5}, $p_{\mathrm{ex}} = 1$), to reduce a
false risk claim from prior $\pR(0)$ to posterior
$\pR \leq p_{\mathrm{clear}}$ when $H_0$ is true and all $n$
tests produce $O_n = 0$ (no spurious false positives), the
required number of tests satisfies
\begin{equation}
\label{eq:n_clear}
N_{\mathrm{clear}} \geq \left\lceil
  \frac{\log \bigl(p_{\mathrm{clear}} (1 - \pR(0)) / (\pR(0)(1 - p_{\mathrm{clear}}))\bigr)}
       {\log \Lzero}
\right\rceil,
\end{equation}
where $\Lzero = [\thetaone \beta + (1 - \thetaone)(1 - \alphafp)]/(1 - \alphafp)$.
Note $\Lzero < 1$ by Assumption~\ref{ass:T1}, so the
right-hand side is positive.  When $\alphafp \ll 1$,
$\Lzero \approx \thetaone \beta + (1 - \thetaone)$ and the
formula reduces approximately to
$$
N_{\mathrm{clear}} \approx \left\lceil
  \frac{\logit p_{\mathrm{clear}} - \logit \pR(0)}
       {\log(\thetaone \beta + 1 - \thetaone)}
\right\rceil.
$$
\end{corollary}

\begin{proof}[Proof sketch]
From~\eqref{eq:update_neg} applied $n$ times,
$\logit \pR(t_n) = \logit \pR(0) + n \log \Lzero$.  Setting
$\logit \pR(t_n) \leq \logit p_{\mathrm{clear}}$ and solving
for $n$ yields the bound.
\end{proof}

\begin{remark}[Concrete numbers]
\label{rem:concrete_numbers}
For $\thetaone = 1$, $\beta = 0.1$, $\alphafp = 0.01$,
$\pR(0) = 0.5$, and $p_{\mathrm{clear}} = 0.01$:
$\Lzero = (1 \cdot 0.1 + 0 \cdot 0.99) / 0.99 = 0.101$, and
$$
N_{\mathrm{clear}} = \left\lceil
  \frac{\log\bigl(0.01 \cdot 0.5 / (0.5 \cdot 0.99)\bigr)}{\log 0.101}
\right\rceil = \lceil 2.006 \rceil = 3 \text{ tests}.
$$
Three controlled relaxation tests suffice to clear a false risk
claim from a neutral prior to $1\%$ residual risk, given $90\%$
detection sensitivity and $1\%$ false-positive rate.  The
correction from including $\alphafp$ adds approximately one
test compared with the naive $\beta^N$ estimate.  This
illustrates the power of direct observation: the verification
asymmetry that makes the Wamura problem seem intractable is
broken by a small number of well-designed tests.
\end{remark}

\begin{corollary}[EWMA-Risk-Trust Convergence]
\label{cor:ewma}
When test outcomes are fed into the risk-trust EWMA
of~\citep[Definition~4]{lasser2026scm} rather than a pure
Bayesian update, the $L^2$ convergence bound below is
regime-independent: it depends on the EWMA forgetting factor
$\alpha_{\mathrm{risk}}$ and the dependency constant
$C_{\mathrm{dep}}$, not on the per-observation likelihood
form, so it holds under both the full-exercise and
partial-exercise regimes of Assumption~\ref{ass:T5}.  The
risk-trust variance $\sigma_j^{\mathrm{risk},2}(t_n)$
converges to $\mu_j^{\mathrm{risk}}$ in $L^2$ at the rate
established in~\citep[Proposition~2]{lasser2026scm}, with the
variance bound
$$
\limsup_{t \to \infty}
  \Var\bigl[\sigma_j^{\mathrm{risk},2}(t)\bigr]
  \leq \frac{1 - \alpha_{\mathrm{risk}}}{1 + \alpha_{\mathrm{risk}}}
       \, C_{\mathrm{dep}} \, \nu_j^{\mathrm{risk}}.
$$
Controlled-relaxation observations have lower variance than
passive observations (test conditions are controlled), so
$C_{\mathrm{dep}}$ may be smaller than in the passive case,
tightening the variance bound.
\end{corollary}
